\documentclass[a4paper]{article}
% Español (México) con XeLaTeX: fontspec para fuentes Unicode, polyglossia
% para la división silábica y los nombres del documento en la variante mexicana.
\usepackage{fontspec}
\usepackage{polyglossia}
\setdefaultlanguage[variant=mexican]{spanish}
% Los nombres chinos van como «pinyin (original)»: xeCJK da a los caracteres
% Han su propia fuente; sin ella XeLaTeX los omite con un aviso.
% FandolSong no cubre algunos caracteres de nombres (U+73FA); WenQuanYi Zen Hei sí.
\usepackage[AutoFallBack=true]{xeCJK}
\setCJKmainfont{FandolSong-Regular.otf}
\setCJKfallbackfamilyfont{\CJKrmdefault}{WenQuanYi Zen Hei}
% polyglossia no traduce los nombres de listings.
\AtBeginDocument{\providecommand{\lstlistingname}{}\renewcommand{\lstlistingname}{Listado}%
  \providecommand{\lstlistlistingname}{}\renewcommand{\lstlistlistingname}{Índice de listados}}
\usepackage{amsmath, amssymb}
\usepackage{graphicx}
\usepackage[margin=2.35cm]{geometry}
\usepackage[most]{tcolorbox}
\usepackage{booktabs}
\usepackage{tabularx}
\usepackage{array}
\usepackage{float}
\usepackage{xcolor}
\usepackage{hyperref}
\usepackage{fancyhdr}
\setCJKmainfont[AutoFakeBold=2.4]{AR PL UMing CN}
\setCJKsansfont[AutoFakeBold=2.4]{Droid Sans Fallback}
\setCJKmonofont[AutoFakeBold=2.4]{Droid Sans Fallback}
\hypersetup{colorlinks=true, linkcolor=blue!55!black, urlcolor=blue!60!black}
\graphicspath{{./}{figures/}}
\setlength{\parskip}{0.55em}
\setlength{\parindent}{2em}
\setlength{\emergencystretch}{3em}
\renewcommand{\arraystretch}{1.18}
\newtcolorbox{knowledgebox}[1]{enhanced, colback=blue!5!white, colframe=blue!70!black, colbacktitle=blue!70!black, coltitle=white, fonttitle=\bfseries, title={#1}, sharp corners, breakable}
\newtcolorbox{importantbox}[1]{enhanced, colback=yellow!10!white, colframe=yellow!75!black, colbacktitle=yellow!75!black, coltitle=black, fonttitle=\bfseries, title={#1}, sharp corners, breakable}
\newtcolorbox{warningbox}[1]{enhanced, colback=red!5!white, colframe=red!70!black, colbacktitle=red!70!black, coltitle=white, fonttitle=\bfseries, title={#1}, sharp corners, breakable}
\pagestyle{fancy}
\fancyhf{}
\fancyfoot[C]{\thepage}
\fancyfoot[R]{\footnotesize\color{gray}\href{https://github.com/hqhq1025/ai-course-notes}{AI Course Notes}}
\renewcommand{\headrulewidth}{0pt}
\renewcommand{\footrulewidth}{0pt}
\newcommand{\videourl}{https://www.youtube.com/watch?v=qW-kgogQwJc}
\newcommand{\repourl}{https://github.com/hqhq1025/ai-course-notes}

\begin{document}
\begin{titlepage}
\centering
{\Large Notas didácticas de una entrevista a fondo\par}
\vspace{1.0cm}
{\huge\bfseries Rokid, los lentes inteligentes y el bosque oscuro del emprendimiento en hardware}\par
\vspace{0.5cm}
{\Large Zhang Xiaojun conversa con Zhu Mingming (祝铭明) Misa, fundador de Rokid\par}
\vspace{0.35cm}
{\large Elaborado a partir del video público, la transcripción, la estructura de capítulos de Zhang Xiaojun Podcast y diagramas conceptuales propios\par}
\vspace{0.3cm}
{\large 2025-06-16\par}
\vspace{0.85cm}
\includegraphics[width=0.88\textwidth,height=0.42\textheight,keepaspectratio]{cover.jpg}\par
\vfill
\begin{tcolorbox}[width=0.92\textwidth, colback=black!2!white, colframe=black!55, sharp corners]
\textbf{Título del video}: 104. Conversación con Zhu Mingming de Rokid: Wu Ma (吴妈), Alibaba y el undécimo año en el bosque oscuro del emprendimiento en hardware\par
\textbf{Canal del video}: Zhang Xiaojun Podcast\par
\textbf{Duración del video}: 02:08:56\par
\textbf{Enlace del video}: \href{\videourl}{\nolinkurl{\videourl}}\par
\textbf{Notas sobre el material}: YouTube no tiene subtítulos disponibles; el cuerpo del texto se elaboró a partir de la transcripción por capítulos con faster-whisper large-v3 CPU int8, la descripción del video, la portada y figuras didácticas propias. Las tomas fijas de la entrevista no se reutilizan en el cuerpo del texto.\par
\textbf{Página del proyecto}: \href{\repourl}{\nolinkurl{\repourl}}\par
\end{tcolorbox}
\end{titlepage}

\begingroup
\small
\setlength{\parskip}{0pt}
\renewcommand{\baselinestretch}{0.92}\selectfont
\tableofcontents
\endgroup
\newpage
\section{Introducción: cómo las capacidades de software de la IA se extienden al hardware}

Esta sección establece primero los objetivos de aprendizaje de todo el episodio. Este episodio sobre Rokid no es una simple historia oral de emprendimiento, sino un caso de estudio sobre «el punto de acceso del hardware de IA de próxima generación». Zhu Mingming (祝铭明) pasó por cuatro ciclos, el internet móvil, las bocinas con IA, los lentes de AR y los modelos grandes: su primera empresa fue adquirida por Alibaba, participó en M Lab dentro de Alibaba y emprendió por segunda vez con Rokid. Al leer este episodio, lo importante no es recordar en qué año hubo una ronda de financiamiento, sino entender cómo un emprendedor de hardware evalúa el punto de acceso, la plataforma, la cadena de suministro, la organización y la competencia con los gigantes.

Zhang Xiaojun (张小珺) dice al inicio que, a medida que las capacidades de software de la IA se extienden al hardware, además de la inteligencia corporizada, los lentes inteligentes podrían ser otra industria beneficiada. Este juicio sitúa el EP104 dentro de la serie de IA/internet de Zhang Xiaojun: no es una entrevista de electrónica de consumo pura, sino una discusión sobre cómo las capacidades de la IA pasan de la nube, el teléfono celular y los productos de voz a un dispositivo personal always-on.

\begin{importantbox}{Tesis central de este episodio}
La oportunidad de los lentes inteligentes no consiste solo en «meter la IA en unos lentes», sino en redefinir el punto de acceso personal a la información. El juicio de Rokid es este: la IA necesita un soporte disponible en cualquier momento y lugar, que entregue la información de forma directa y que tenga capacidad de visualización; el teléfono celular seguirá existiendo, pero las interacciones fragmentadas podrían trasladarse poco a poco a los lentes.
\end{importantbox}

\begin{figure}[H]
\centering
\includegraphics[width=0.96\textwidth]{ai-to-hardware-spillover.png}
\caption{Las capacidades de software de la IA se extienden al hardware: del modelo en la nube a la inteligencia portátil y la inteligencia corporizada. Diagrama conceptual propio, elaborado a partir de la conversación en 00:01:12--00:02:00 y 01:11:52--01:13:05.}
\end{figure}

\begin{knowledgebox}{Lectura de la figura: la inteligencia corporizada y la inteligencia portátil son dos rutas hacia el hardware}
En la figura, AI Model se extiende hacia Device y se divide en dos categorías: Wearable y Embodied. La inteligencia corporizada pone el énfasis en que los robots actúen en el mundo físico; la inteligencia portátil, en que las personas puedan recurrir a la IA en cualquier momento. Los lentes inteligentes pertenecen a la segunda: lo decisivo en ellos es la eficiencia como punto de acceso, la visualización y la comodidad al llevarlos puestos, no la ejecución mecánica.
\end{knowledgebox}

\begin{knowledgebox}{Nota sobre la estrategia visual}
Este video es una entrevista con encuadre fijo, sin slides, pizarrón ni demostraciones de producto. El texto solo usa la portada para identificar la fuente; todas las imágenes del cuerpo son diagramas conceptuales propios que explican la cronología del emprendimiento, el giro hacia la IA y la AR, el «bosque oscuro» del hardware, las etapas de los lentes inteligentes y las diferencias en la definición de producto entre China y Estados Unidos.
\end{knowledgebox}

\subsection{Ruta de lectura: la historia del emprendimiento como lección sobre cómo evaluar puntos de acceso}

Esta sección propone además una forma de leer. La entrevista contiene una gran cantidad de nombres de personas, rondas de financiamiento, productos y años; si se memorizan uno por uno, es fácil leerla como una «historia del fundador». Una lectura más valiosa es devolver cada episodio de la historia a un marco de juicio: cuando un nuevo hardware busca convertirse en punto de acceso, qué condiciones debe cumplir exactamente; y cuando un producto solo está aprovechando una tendencia, qué señales indican que es un producto y no una plataforma.

\begin{knowledgebox}{Nota de clase: este episodio no es una evaluación de especificaciones de lentes}
Aquello a lo que el ponente vuelve una y otra vez no es una especificación concreta, sino el juicio abstracto sobre el «punto de acceso»: tiempo de uso, amplitud de entrada y salida, frecuencia de las interacciones fragmentadas, naturalidad al llevarlos puestos, capacidad de visualización, posición en el ecosistema y capacidad de entrega de la cadena de suministro. Las especificaciones solo están al servicio de estos juicios; no pueden, a la inversa, sustituirlos.
\end{knowledgebox}

\noindent\begin{tabularx}{\textwidth}{p{0.18\textwidth}p{0.36\textwidth}X}
\toprule
Pregunta de lectura & Material en la entrevista & Juicio que se debe formar \\
\midrule
¿De dónde surge el punto de acceso? & iPhone, sistemas operativos móviles, bocinas con IA, lentes inteligentes & El cambio de punto de acceso suele ocurrir cuando la forma del hardware y las capacidades del software cambian al mismo tiempo. \\
¿Cuándo se consolida una plataforma? & Dos horas de uso diario promedio, desarrolladores, usuarios empresariales, ecosistema de contenido & Una plataforma no existe porque la empresa lo proclame, sino que se forma por la combinación de la frecuencia de uso y la oferta del ecosistema. \\
¿Por qué hace falta visualización? & Traducción, avisos, consulta rápida de información, resultados intermedios de baja confianza & La visualización convierte la salida de la IA, de voz en serie, en una interfaz visible y corregible. \\
¿Cómo sobrevive una empresa emergente? & Los cuatro «no», las ventanas de oportunidad, «acumular grano y levantar murallas altas», hacer muchos aliados & En la etapa temprana, la ventana se aprovecha con velocidad; la competencia a largo plazo depende del producto, el ecosistema y la paciencia del capital. \\
\bottomrule
\end{tabularx}

\subsection{Resumen de la sección}

El hilo conductor del EP104 es este: cuando las capacidades de la IA maduran, los puntos de acceso de hardware se reordenan. El valor de Rokid como caso de estudio está en que, con el costo de más de diez años de emprendimiento en hardware, explica por qué los lentes inteligentes podrían convertirse en el punto de acceso de la IA portátil. Las secciones siguientes explicarán primero de dónde proviene este juicio sobre el punto de acceso y, después, por qué se abandonaron las bocinas y se persistió en los lentes, cómo llegó la competencia de los gigantes y por qué al emprendimiento en hardware se le llama un «bosque oscuro».
\section{De la primera empresa a Alibaba: los antecedentes de los puntos de entrada del internet móvil}

Esta sección examina primero el primer emprendimiento de Zhu Mingming (祝铭明), porque explica por qué Rokid siempre ha tenido la mira puesta en el «punto de entrada de la siguiente generación». Después del lanzamiento del iPhone en 2007, él concluyó que los teléfonos inteligentes y el internet móvil se convertirían en una nueva plataforma, así que fundó una empresa de sistemas operativos móviles con el propósito de trasladar un sistema de software similar al del iPhone a distintas plataformas de hardware. Más tarde, Alibaba se preparaba para entrar al internet móvil y necesitaba un punto de apoyo básico; al final adquirió esa empresa por unos 10 millones de dólares.

Lo esencial de esta historia no es «venderse a Alibaba», sino tres lecciones. Primera: el cambio de punto de entrada suele ocurrir en el momento en que el hardware y el sistema de software cambian al mismo tiempo. Segunda: una empresa emergente puede llamar la atención de una gran empresa justo en su momento más crítico. Tercera: las oportunidades a escala de plataforma exigen que el fundador soporte una incertidumbre altísima en las etapas iniciales. Zhu Mingming recuerda que, en el peor momento, la empresa tenía menos de 4,000 yuanes en cuenta, la nómina debía pagarse en dos semanas e incluso se preparaba para que los empleados se llevaran a casa las computadoras como pago de sus salarios.

\begin{figure}[H]
\centering
\includegraphics[width=0.96\textwidth]{rokid-founder-timeline.png}
\caption{Línea de tiempo del emprendimiento de Rokid: de la adquisición por Alibaba y M Lab a once años de hardware de IA/AR. Diagrama conceptual de elaboración propia, organizado a partir de la conversación entre 00:02:36 y 00:31:41.}
\end{figure}

\begin{knowledgebox}{Lectura de la figura: el primer emprendimiento no fue un episodio aislado, sino un entrenamiento para identificar puntos de entrada}
La figura va de First Startup a Alibaba y, después, a M Lab y Rokid, y muestra un mismo tipo de juicio: cuando surge una nueva plataforma de cómputo, el sistema de base, el hardware de entrada y el ecosistema de aplicaciones se reorganizan por completo. Zhu Mingming pasó del OS móvil a Alibaba y después salió para dedicarse a la IA/AR; en esencia, siempre ha buscado el punto de entrada de la siguiente generación.
\end{knowledgebox}

\subsection{Por qué importan los momentos de crisis}

Esta sección aparta para examinarla por separado la historia inicial de «menos de cuatro mil yuanes en cuenta». No es un pasaje motivacional, sino un entrenamiento en flujo de efectivo para quien emprende en torno a un punto de entrada. Zhu Mingming recuerda que en ese momento la nómina debía pagarse en dos semanas e incluso se preparaba para que los empleados se llevaran a casa las computadoras como pago de sus salarios; justo en esa etapa, Ma Yun (马云), Cai Chongxin (蔡崇信) y el doctor Wang Jian (王坚) abrieron el camino hacia la adquisición por parte de Alibaba. Que una empresa emergente llame la atención no siempre ocurre en su momento más brillante; puede ocurrir, por el contrario, cuando solo unas pocas personas entienden su dirección y el flujo de efectivo ya está cerca del límite.

\begin{warningbox}{Experiencia práctica: una dirección correcta no garantiza un flujo de efectivo seguro}
Aunque un proyecto emprendedor apueste correctamente por el punto de entrada de la siguiente generación, puede agotar su efectivo antes de que lleguen la comercialización, el financiamiento o los grandes clientes. Tanto la primera empresa de OS móvil mencionada en la entrevista como la posterior Rokid lo muestran: las direcciones de vanguardia requieren paciencia estratégica, pero la empresa sigue enfrentando cada día la nómina, los proveedores, el inventario y la confianza del equipo.
\end{warningbox}

\noindent\begin{tabularx}{\textwidth}{p{0.20\textwidth}p{0.34\textwidth}X}
\toprule
Suceso & Historia aparente & Implicación didáctica \\
\midrule
Lanzamiento del iPhone en 2007 & El fundador ve la dirección de los teléfonos inteligentes & Cuando aparece un nuevo punto de entrada, el hardware y el sistema de software anteriores pierden firmeza. \\
Desarrollo de un OS móvil & Trasladar un sistema de software al estilo del iPhone a varios modelos de teléfono & Las oportunidades de plataforma suelen manifestarse primero como problemas del sistema de base y de compatibilidad. \\
Adquisición por Alibaba & Operación en acciones por 10 millones de dólares & Cuando un gigante quiere entrar a una nueva plataforma, necesita capacidades básicas y un equipo. \\
Crisis de flujo de efectivo & No alcanza el dinero en cuenta para pagar la nómina & Por ambicioso que sea el juicio sobre el punto de entrada, la empresa tiene que superar el umbral de supervivencia. \\
\bottomrule
\end{tabularx}

\subsection{Alibaba, Wu Ma (吴妈) y M Lab}

Esta sección explica cómo la experiencia en Alibaba influyó en Rokid. Tras incorporarse a Alibaba, Zhu Mingming participó al principio en negocios relacionados con el internet móvil; después, Wu Yongming (吴泳铭) propuso trabajar en IA y se fundó M Lab. En ese entonces el aprendizaje profundo apenas comenzaba, y el reconocimiento de voz, la búsqueda por fotografía, el AIoT y el hardware inteligente estaban en una fase temprana de exploración. Alibaba quería hacer algo fundamental: plataformas, sistemas y capacidades básicas, no solo una aplicación puntual.

La experiencia de M Lab dejó en Rokid dos influencias duraderas. Primera: la IA no es una función aislada, sino que busca un punto de entrada en el hardware. Segunda: una gran empresa puede organizar recursos, pero en ciertas direcciones de hardware de vanguardia una empresa emergente es más adecuada para adelantarse con alto riesgo y alta concentración. Wu Yongming fue después inversionista ángel de Rokid, lo que muestra que esta relación no es solo un antecedente de trayectoria, sino que también influyó en el financiamiento y en las decisiones posteriores.

\begin{knowledgebox}{Términos clave: punto de entrada, plataforma, producto}
\noindent\begin{tabularx}{\textwidth}{p{0.18\textwidth}p{0.34\textwidth}X}
\toprule
Concepto & Significado & Juicio en este episodio \\
\midrule
Producto & Resuelve una necesidad concreta; sus ingresos provienen de la venta de unidades o de servicios & Una bocina con IA se acerca más a un producto; no necesariamente es una plataforma. \\
Punto de entrada & Canal predeterminado por el que el usuario entra con frecuencia al mundo digital & El teléfono es el punto de entrada de la generación anterior; los lentes podrían asumir el punto de entrada fragmentado. \\
Plataforma & Capa de base capaz de sostener desarrolladores, aplicaciones y un ecosistema & Los lentes inteligentes solo podrían convertirse en plataforma si alcanzan suficiente tiempo de uso y suficientes desarrolladores. \\
\bottomrule
\end{tabularx}
\end{knowledgebox}

\begin{importantbox}{El docente enfatiza: lo fundamental cambia los límites de la organización}
En aquel momento Alibaba absorbió al equipo del OS móvil porque el internet móvil requería un punto de apoyo en la capa de base; Wu Yongming impulsó la IA y M Lab porque la IA no podía tratarse solo como un complemento funcional. Que Rokid haya insistido después en el OS, la plataforma y el ecosistema proviene de la misma experiencia: la tecnología que realmente cambia el punto de entrada termina por tocar el sistema de base.
\end{importantbox}

\subsection{Resumen de la sección}

Los antecedentes de Zhu Mingming en Alibaba desarrollaron dos capacidades: juzgar cuándo cambia el punto de entrada de una plataforma y elegir, entre la gran empresa y la empresa emergente, la forma organizativa que le conviene. Esto sentó las bases de las dos líneas posteriores, IA y AR.
\section{El giro de 2019: de las bocinas con IA a los lentes de AR}

La sección anterior explicó de dónde surge el juicio de Rokid sobre el punto de entrada; esta sección aborda el giro estratégico más importante. Desde el primer día, Rokid consideró que la IA y la AR terminarían siendo una sola cosa, pero el hardware de AR de los primeros años era demasiado pesado y extraño, y no podía convertirse en producto. Por eso, la empresa empezó con hardware de consumo con IA y bocinas inteligentes para desarrollar sus capacidades en voz, cadena de suministro, producción, ventas y retroalimentación de usuarios. En 2019, cuando las grandes empresas entraron al mercado de las bocinas inteligentes con subsidios y con una lógica de plataforma, Rokid concluyó que la bocina no podía convertirse en un verdadero punto de entrada y, en una semana, pasó del sector de la IA al de la AR.

Este giro fue muy doloroso, pero no fue un cambio repentino de dirección. Zhu Mingming (祝铭明) subraya que internamente siempre avanzaron sobre dos piernas en paralelo, la IA y la AR; la diferencia era que antes la línea de productos de IA podía «financiar la guerra con la propia guerra». Cuando la competencia en bocinas se convirtió en un gasto de capital poco sano, seguir invirtiendo dejó de ser razonable. Lo que Rokid quería librar era la guerra del futuro, no consumir efectivo en un campo de batalla de plataformas que no le pertenecía.

\begin{figure}[H]
\centering
\includegraphics[width=0.94\textwidth]{ar-pivot-decision.png}
\caption{2019: el giro de las bocinas con IA hacia la AR, con el cambio de sector en una semana; el costo fue la reestructuración de la organización. Diagrama conceptual propio, elaborado a partir de la conversación entre 00:31:41 y 00:59:17.}
\end{figure}

\begin{knowledgebox}{Lectura de la figura: no es abandonar la IA, sino cambiar de soporte}
A la izquierda, la bocina con IA es un producto de voz relativamente maduro que permite formar al equipo, pero sus canales de entrada y salida son estrechos, sus escenarios de uso no son fragmentados y su carácter de plataforma es débil. A la derecha, los lentes de AR son más difíciles, pero están más cerca de un punto de entrada de IA que se lleva encima. El giro de Rokid consiste en pasar de un producto viable a corto plazo a un punto de entrada más grande a largo plazo.
\end{knowledgebox}

\subsection{Dos piernas en paralelo: la línea de productos de IA y el laboratorio de AR}

Esta subsección añade un detalle que se malinterpreta con facilidad. Zhu Mingming explica que Rokid no pasó de pronto de la IA a la AR, sino que desde el principio consideró que la IA y la AR convergerían en una sola cosa. Los primeros dispositivos de AR eran demasiado pesados y parecían montajes de laboratorio, por lo que no podían ser productos para el gran público; en cambio, las tecnologías de voz como NLP, ASR y TTS ya podían convertirse en producto. Así, el hardware de IA asumió primero dos tareas: mantener a la empresa en contacto con el mercado y desarrollar sus capacidades de cadena de suministro de hardware, producción, ventas y retroalimentación de usuarios.

Ese es el verdadero sentido de «financiar la guerra con la propia guerra». La bocina no era la plataforma final, pero permitía que un equipo con origen en software y sistemas aprendiera a hacer hardware. Zhu Mingming menciona que de 2014 a 2016 en realidad estuvieron pagando el precio de aprender hardware: la definición del producto, el precio, la cadena de suministro y la aceptación en el mercado de consumo tenían que pasar la prueba del mercado. El primer robot de compañía de Rokid, Alien, era muy atractivo, pero estaba demasiado adelantado a su tiempo; al final se simplificó hasta convertirse en una bocina, con lo que también bajó la barrera de entrada, y cuando los gigantes llegaron con subsidios, que una empresa emergente siguiera quemando dinero dejó de ser razonable.

\noindent\begin{tabularx}{\textwidth}{p{0.18\textwidth}p{0.32\textwidth}X}
\toprule
Ruta & Función en ese momento & Problema o valor posterior \\
\midrule
Robot de compañía Alien & Reunir cámara, expresiones, voz y compañía en un solo producto & Concepto adelantado a su tiempo, pero con precio y aceptación del mercado de consumo insuficientes. \\
Bocina con IA & Formar capacidades en productos de voz, producción, ventas y cadena de suministro & Al bajar la barrera de entrada, cayó en una competencia de subsidios; carácter de plataforma insuficiente. \\
AR Lab & Investigación previa a largo plazo en guías de onda ópticas, micropantallas, formación de imagen y otras tecnologías & Hizo que el giro de 2019 no partiera de cero. \\
Lentes inteligentes & Volver a sostener la hipótesis del punto de entrada de IA que se lleva encima & Más difíciles, pero con un techo de plataforma más alto. \\
\bottomrule
\end{tabularx}

\begin{knowledgebox}{Nota de clase: empezar por un producto no significa renunciar al ideal de plataforma}
Rokid hizo primero bocinas porque la IA de voz podía llevarse a la práctica antes; más tarde abandonó las bocinas porque no podían sostener la hipótesis de plataforma. Este orden muestra que los productos de cada etapa de una empresa emergente pueden servir para construir capacidades a largo plazo, pero no deben, a la inversa, secuestrar la dirección de largo plazo.
\end{knowledgebox}

\subsection{Por qué la bocina no es un punto de entrada}

Esta subsección explica con claridad la lógica del fracaso de la bocina. Zhu Mingming sostiene que un producto de entrada necesita canales de entrada y salida suficientemente amplios, suficientes escenarios de uso y una frecuencia de uso fragmentado suficientemente alta. Los canales de la bocina son estrechos: sobre todo lectura, música y preguntas sencillas; sus escenarios también se limitan a la sala y la recámara, y el usuario pasa la mayor parte de sus horas de vigilia fuera de esos escenarios fijos. Puede ser un buen producto, pero difícilmente se convierte en un punto de entrada del nivel del teléfono celular.

Las grandes empresas trataron la bocina inteligente como una plataforma: subsidios, gasto de capital y disputa por el ecosistema, lo cual perjudica a las empresas emergentes. Rokid no consideraba la bocina una plataforma, así que se retiró con decisión. Este juicio es importante: una empresa emergente no puede fijarse solo en si «el mercado está caliente»; también debe juzgar si ese entusiasmo concuerda con su propia hipótesis de plataforma.

\begin{warningbox}{La ilusión de plataforma}
Que un producto tenga IA, voz y hardware no significa que sea una plataforma. Una plataforma requiere un punto de entrada de uso frecuente, un canal de información suficientemente amplio, un ecosistema sostenible y hábitos de uso predeterminado por parte de los usuarios. Sin estas condiciones, gastar capital puede limitarse a amplificar un producto común.
\end{warningbox}

\begin{knowledgebox}{Lista de criterios para evaluar un punto de entrada}
\noindent\begin{tabularx}{\textwidth}{p{0.22\textwidth}p{0.32\textwidth}X}
\toprule
Criterio & Problema de la bocina & Posibilidades de los lentes \\
\midrule
Amplitud de entrada y salida & Depende sobre todo de la voz, y la salida también tiende a ser secuencial & Voz, visión, pantalla y percepción del entorno pueden sumarse. \\
Frecuencia de uso fragmentado & Fija en escenarios limitados como la recámara o la sala & Se llevan puestos y pueden cubrir el trayecto, las tiendas, las reuniones y los viajes. \\
Ruta para sustituir al teléfono & Para muchas tareas, el teléfono ya es bastante cómodo & Tareas breves como ver la hora, escanear códigos, traducir o recibir notificaciones pueden acortar la ruta. \\
Espacio de ecosistema & Los tipos de contenido se concentran en música, lectura y preguntas & Si el tiempo de uso y la pantalla se consolidan, el ecosistema de aplicaciones es más amplio. \\
Ventaja de la empresa emergente & Los gigantes pueden subsidiar y distribuir a gran escala & Los procesos de fabricación de frontera y la definición del producto aún no están estandarizados. \\
\bottomrule
\end{tabularx}
\end{knowledgebox}

\subsection{Despidos y reestructuración de la organización}

Esta subsección pasa del juicio sobre el producto al costo organizacional. El costo del giro estratégico fue la reestructuración de la organización. En ese momento Rokid tenía dos edificios, uno dedicado a las bocinas y otro a la AR; tras la transformación, el edificio de las bocinas quedó prácticamente vacío y se despidió a más de la mitad del personal. Zhu Mingming subraya que lo doloroso no fue el juicio sobre la dirección en sí, sino que muchos empleados excelentes tuvieron que irse por el cambio de dirección. La empresa ofreció entonces una indemnización N+1 y dejó abierta la posibilidad de que regresaran con prioridad una vez que la AR funcionara bien.

Este episodio es muy instructivo para las empresas emergentes de hardware. En un producto de software el giro puede ser más ligero; en una línea de productos de hardware, el giro suele afectar la cadena de suministro, las habilidades del equipo, el inventario, los planes de producción y el flujo de efectivo. Rokid pudo dar el giro porque aún tenía dinero en caja, sus inversionistas entendían la dirección de largo plazo y el AR Lab llevaba ya varios años funcionando. Si hubiera esperado a agotar por completo el efectivo, la ventana para transformarse habría desaparecido.

\begin{importantbox}{La ventana de transformación}
Una empresa de hardware que emprende un gran giro debe iniciarlo cuando todavía tiene efectivo suficiente, el equipo aún conserva sus capacidades y la siguiente ruta ya cuenta con investigación previa. Si espera a que la línea de productos fracase por completo, la organización, la cadena de suministro y la confianza de los inversionistas suelen ser ya insuficientes.
\end{importantbox}

\subsection{El primer PMF durante la pandemia}

A continuación se examinan el azar y la necesidad después de la transformación. Rokid hizo el ajuste en octubre de 2019 y, a mediados de enero de 2020, celebró un evento de captación de distribuidores para sus lentes; pocos días después estalló la pandemia y se restringió la movilidad presencial. Como los lentes podían combinarse con sensores infrarrojos para medir la temperatura sin contacto, la primera generación de productos para clientes empresariales, más pesada, se industrializó rápidamente y se vendió en varios países. Zhu Mingming no lo describe como «gracias a la pandemia», pero reconoce que, objetivamente, aceleró las curvas del producto, de la cadena de suministro y de los ingresos.

Lo más valioso de este episodio es distinguir entre PMF y suerte. La pandemia fue un acontecimiento externo que Rokid no podía prever; pero si en 2019 no se hubiera pasado a tiempo a la AR, no hubiera celebrado con anticipación el evento de captación de distribuidores ni hubiera preparado la cadena de suministro, no habría podido aprovechar ese acontecimiento. En las empresas emergentes de hardware, el PMF a menudo no es una pura inspiración de producto, sino la alineación repentina de la ruta, el momento, la cadena de suministro y el escenario.

\begin{knowledgebox}{Experiencia práctica: el azar hay que estar preparado para aprovecharlo}
Zhu Mingming dice que si el evento de captación de distribuidores se hubiera celebrado cinco días después, el resultado podría haber sido completamente distinto. No se trata de alentar una creencia supersticiosa en el momento oportuno, sino de recordar que una empresa de hardware de frontera debe completar suficientes preparativos antes de que se abra la ventana. Cuando surge la oportunidad, ser capaz de entregar suele importar más que ser capaz de contar una historia.
\end{knowledgebox}

\subsection{Resumen de la sección}

El giro de Rokid en 2019 no fue de la IA a algo ajeno a la IA, sino de «productos de IA de voz» a «el punto de entrada de la IA que se lleva encima». Este giro formó el juicio de la empresa sobre el carácter de plataforma, la cadena de suministro y el costo organizacional, y muestra también que una transformación de hardware debe encontrar su ventana entre el financiamiento, la investigación previa, la organización y el momento externo.
\section{Por qué las gafas inteligentes podrían convertirse en un punto de entrada}

La sección anterior explicó por qué Rokid dejó las bocinas inteligentes; esta sección explica por qué apuesta por las gafas. Zhu Mingming (祝铭明) considera que la forma de usar la IA que realmente aporta valor es en cualquier momento y lugar, fragmentada y always on. El teléfono es potente, pero sigue exigiendo un recorrido de GUI: desbloquear, abrir una app y entrar en una función. Si las gafas inteligentes se llevan puestas todo el tiempo, muchas tareas breves pueden volverse «lo que pides es lo que obtienes»: consultar la hora, recibir mensajes, escanear códigos, traducir, reconocer objetos o consultar información de contexto, todo sin sacar el teléfono.

Esto no significa que el teléfono vaya a desaparecer de inmediato, sino que la frecuencia de interacción se redistribuirá. El teléfono podría conservar la comunicación, las tareas pesadas, la entrada compleja y las operaciones en pantalla grande; las gafas absorberían primero gran parte de las interacciones fragmentadas de menos de cinco minutos. El juicio de Zhu Mingming es este: para quienes ya usan lentes, las gafas con IA podrían volverse algo de uso generalizado en un plazo de tres a cinco años; a quienes no quieren usar lentes no hay que intentar convencerlos a la fuerza, sino esperar a que la diferencia de capacidades sea lo bastante evidente y dejar que ellos decidan.

\begin{figure}[H]
\centering
\includegraphics[width=0.96\textwidth]{smart-glasses-stage.png}
\caption{Etapa de las gafas inteligentes: entre BlackBerry y el iPhone 1, el punto de entrada está tomando forma. Diagrama conceptual propio, elaborado a partir de la conversación en 01:05:45--01:13:05.}
\end{figure}

\begin{knowledgebox}{Lectura de la figura: que la etapa sea inmadura no significa que la dirección sea errónea}
La figura va de Feature a BlackBerry, iPhone 1, Platform y Mass Market, y expresa que las gafas inteligentes todavía están en la etapa temprana en la que se forma un punto de entrada. Los productos actuales no son necesariamente maduros, pero si se confirman el tiempo de uso, los desarrolladores y la migración de tareas fragmentadas, tienen la oportunidad de pasar de hardware de función a plataforma.
\end{knowledgebox}

\subsection{Acceso directo a la información: reducir al mínimo el costo del sistema}

Esta sección descompone la «comodidad» en costos de interacción analizables. Zhu Mingming pone como ejemplos escanear un código, consultar el clima, ver la hora y leer mensajes: en el teléfono, el recorrido suele incluir sacar el dispositivo, desbloquearlo, encontrar la app, entrar en la función y confirmar la operación; con unas gafas que se llevan puestas todo el tiempo, basta con invocarlas por voz, confirmar visualmente o tocarlas ligeramente para ver el resultado. Unos segundos de diferencia importan poco en una sola tarea, pero cuando todas las tareas son acciones fragmentadas de menos de cinco minutos, esos segundos se acumulan hasta provocar la migración del punto de entrada.

\[
T_{\text{task}} = T_{\text{wake}} + T_{\text{path}} + T_{\text{input}} + T_{\text{output}} + T_{\text{confirm}}
\]

Aquí, $T_{\text{task}}$ es el tiempo total para completar una tarea breve; $T_{\text{wake}}$ es el tiempo para sacar o activar el dispositivo; $T_{\text{path}}$ es el costo del recorrido desde el punto de entrada del sistema hasta la función concreta; $T_{\text{input}}$ es el tiempo para expresar la intención; $T_{\text{output}}$ es el tiempo para recibir el resultado; $T_{\text{confirm}}$ es el tiempo para confirmar y corregir errores. Lo esencial de la vía de las gafas no es llevar cada término a cero, sino reducir mucho $T_{\text{wake}}$ y $T_{\text{path}}$, y usar la pantalla para disminuir la incertidumbre de $T_{\text{output}}$ y $T_{\text{confirm}}$.

\begin{importantbox}{El punto de entrada no depende de la eficiencia de una sola vez, sino de la frecuencia multiplicada por la eficiencia}
Si una acción ocurre una sola vez al día, ahorrar cinco segundos no tiene importancia; si las notificaciones, la hora, el escaneo de códigos, la traducción, el reconocimiento y la búsqueda ocurren decenas de veces al día, el costo del recorrido se convierte en un criterio para elegir el punto de entrada. El juicio de Rokid es que las gafas inteligentes primero sustituirán una gran cantidad de tareas ligeras y fragmentadas, no las tareas pesadas del teléfono de inmediato.
\end{importantbox}

\subsection{Por qué importa la pantalla}

Una de las diferencias centrales en la definición de producto de las gafas inteligentes entre China y Estados Unidos es si necesitan pantalla. Rokid considera que, si las gafas son un punto de entrada a la IA, la pantalla es indispensable. Tomemos la traducción: sin pantalla, el sistema debe esperar a tener suficiente confianza antes de leer la traducción en voz alta, y en la conversación aparecen silencios de varios segundos hasta diez segundos; con pantalla, puede mostrar mientras escucha y corregir sobre la marcha, y el usuario tolera resultados intermedios que al principio no son del todo precisos. La pantalla no está para lucirse, sino para que la salida de la IA pase de ser solo voz a una interfaz que se puede leer de un vistazo, corregir y usar en paralelo.

\begin{knowledgebox}{Términos clave: AUI y GUI}
\noindent\begin{tabularx}{\textwidth}{p{0.18\textwidth}p{0.34\textwidth}X}
\toprule
Modo de interacción & Ventajas & Limitaciones \\
\midrule
GUI & Visible, editable y fácil de leer de un vistazo; adecuada para tareas complejas & Requiere recorridos y operación de la interfaz; el costo del sistema es alto. \\
AUI & Acceso directo por voz; adecuada para tareas breves y uso sin manos & La salida es secuencial; no es adecuada para contenido largo ni para resultados intermedios de baja confianza. \\
Pantalla en gafas + IA & Combina entrada por voz y salida visual & Exige que la comodidad al llevarlas puestas, la autonomía de la batería, la pantalla y la privacidad superen juntas el umbral. \\
\bottomrule
\end{tabularx}
\end{knowledgebox}

\begin{warningbox}{No entender la pantalla como «una pantalla más»}
El verdadero valor de la pantalla no es trasladar la del teléfono frente a los ojos, sino cambiar la estructura temporal de la salida de la IA. La salida por voz tiene que ser secuencial, completa y de confianza relativamente alta; la salida visual puede mostrarse de forma gradual, admitir correcciones intermedias y permitir que el usuario la lea de un vistazo. La traducción es un caso representativo porque expone a la vez la latencia, la confianza y la posibilidad de corregir errores.
\end{warningbox}

\noindent\begin{tabularx}{\textwidth}{p{0.18\textwidth}p{0.34\textwidth}X}
\toprule
Escenario & Recorrido sin pantalla & Recorrido con pantalla \\
\midrule
Traducción & Se lee en voz alta solo cuando el modelo tiene suficiente confianza; la conversación entre ambas partes se interrumpe & Puede mostrar primero una traducción aproximada y corregirla según el contexto; el usuario la lee de un vistazo. \\
Notificaciones & El aviso por voz interrumpe la situación, y revisar el teléfono alarga el recorrido & Aviso breve frente a los ojos; el usuario puede ignorarlo o atenderlo. \\
Escaneo de códigos/pago & Hay que abrir el teléfono, buscar el punto de entrada y confirmar & Se dirige la mirada al código QR y se confirma por voz o con un toque ligero. \\
Reconocimiento de objetos & Tomar una foto con el teléfono o abrir la cámara & Las gafas forman la consulta directamente a partir de la dirección de la mirada. \\
\bottomrule
\end{tabularx}

\subsection{Límites de uso: no intentar convencer a quien no quiere usar lentes}

Esta sección aborda una objeción realista: mucha gente no quiere usar lentes, e incluso quienes tienen miopía optan por lentes de contacto o cirugía. La respuesta de Zhu Mingming es mesurada: en la etapa actual no hay que intentar convencer a la fuerza a ese grupo. Primero se atiende a quienes ya aceptan lentes con armazón, sumando capacidades de IA a un hábito que ya tienen; cuando las personas de su entorno que usan lentes obtengan una ventaja informativa evidente, quienes no quieren usarlos podrán volver a decidir. La estrategia de producto no es «que todos los usen de inmediato», sino empezar por el grupo para el que resulta más natural, los usuarios nativos.

\begin{knowledgebox}{Experiencia práctica: la curva de adopción empieza por el grupo con menos resistencia}
Los primeros usuarios de un producto de hardware no son necesariamente el grupo más grande, sino el que opone menos resistencia, da la retroalimentación más intensa y está dispuesto a tolerar las imperfecciones iniciales. Para las gafas inteligentes, quienes están dispuestos a usar lentes con armazón, los entusiastas de la tecnología y los usuarios frecuentes de viajes de trabajo, traducción o reuniones son mejores usuarios principales tempranos que quienes rechazan por completo los lentes.
\end{knowledgebox}

\noindent\begin{tabularx}{\textwidth}{p{0.20\textwidth}p{0.34\textwidth}X}
\toprule
Tipo de usuario & Resistencia inicial & Estrategia de producto \\
\midrule
Quienes ya usan lentes & La más baja; solo hay que convertir el uso existente en capacidad de IA & Priorizar comodidad, graduación, peso y uso diario. \\
Quienes están dispuestos a usarlos en algunos escenarios & Media; los escenarios clave deben ser lo bastante sólidos & Escenarios de alto valor como viajes, traducción, reuniones, navegación y escaneo de códigos. \\
Quienes no quieren usar lentes & La más alta; hay resistencia tanto fisiológica como psicológica & No forzar su convencimiento; ampliar el alcance cuando la diferencia de capacidades y el formato hayan madurado. \\
\bottomrule
\end{tabularx}

\subsection{Diferencias de definición de producto entre China y Estados Unidos}

Muchas gafas inteligentes estadounidenses siguen más bien la lógica de las sunglasses: exteriores, redes sociales, fotos y videos para compartir; no necesariamente son always on ni tratan la pantalla como indispensable. Rokid se inclina más por la lógica china de los lentes: muchos usuarios ya usan lentes, y llevarlos puestos durante largos periodos, en interiores y exteriores, resulta más natural; por eso, desde el primer día, la definición de producto se orienta a IA más AR, con énfasis en la pantalla, el asistente y un punto de entrada a la información que se lleva encima. La vía de Meta/Ray-Ban es una herramienta de producción de contenido social; la de Rokid es el consumo personal de información y el asistente de IA.

\begin{figure}[H]
\centering
\includegraphics[width=0.94\textwidth]{china-us-glasses-definition.png}
\caption{Diferencias en la definición de las gafas inteligentes entre China y Estados Unidos: difieren la definición de producto, los canales y los escenarios de uso. Diagrama conceptual propio, elaborado a partir de la conversación en 01:23:38--01:41:35.}
\end{figure}

\begin{knowledgebox}{Lectura de la figura: que ambas parezcan lentes no significa que sean la misma especie}
A la izquierda, la vía china enfatiza a la población con miopía, el uso durante todo el día, la pantalla y el asistente de IA; a la derecha, la vía estadounidense entra con más facilidad desde las sunglasses, la grabación en redes sociales y la moda de marca. Ambas podrían converger con el tiempo, pero en la etapa temprana difieren en la definición de producto, los escenarios de uso y la prioridad de las funciones.
\end{knowledgebox}

\begin{importantbox}{La definición de producto determina la pila tecnológica}
Si el producto es una herramienta de grabación para redes sociales, lo central son la cámara, la marca, la cadena para compartir y los avisos de privacidad; si el producto es un punto de entrada a la IA, lo central son la pantalla, el sistema operativo, la retroalimentación de baja latencia, el consumo de energía, la comodidad al llevarlas puestas y la integración de modelos. Aunque ambas se llamen gafas inteligentes, sus prioridades técnicas pueden ser completamente distintas.
\end{importantbox}

\subsection{Resumen de la sección}

La lógica de punto de entrada de las gafas inteligentes se basa en estar always-on, en el acceso directo a la información y en la migración de tareas fragmentadas. La capacidad de pantalla determina si la salida de la IA puede pasar de voz secuencial a una interacción visible, corregible y capaz de convertirse en plataforma; los límites de uso, a su vez, determinan que el grupo inicial no puede ampliarse sin límite y que hay que empezar por los usuarios con menos resistencia.
\section{Competir con los gigantes: los cuatro «no» y la ventana de tiempo}

La sección anterior trató la definición del producto; esta trata la competencia. Zhu Mingming (祝铭明) cita el resumen de Ma Yun (马云) sobre las cuatro oportunidades que tiene una empresa emergente frente a los gigantes: no lo ven, no lo entienden, no les interesa y no llegan a tiempo. En el caso de los lentes inteligentes, las grandes empresas sin duda lo ven y muy probablemente también lo entienden; la ventana de la empresa emergente proviene sobre todo de los dos o tres años que transcurren entre «no les interesa» y «no llegan a tiempo». La meta de Rokid no es competir en igualdad de condiciones con los gigantes, sino lograr una competencia entre pares, o al menos «sentarse a la mesa de los adultos».

\begin{figure}[H]
\centering
\includegraphics[width=0.96\textwidth]{startup-vs-giants-four-no.png}
\caption{Los cuatro «no» de la empresa emergente frente a los gigantes: no lo ven, no lo entienden, no les interesa, no llegan a tiempo. Diagrama conceptual propio, elaborado a partir de la conversación de 01:19:29--01:23:38.}
\end{figure}

\begin{knowledgebox}{Lectura de la figura: la ventana de Rokid no consiste en que los gigantes no lo sepan, sino en que todavía no se dedican de lleno}
En la figura, «no llegan a tiempo» es el estado que la empresa emergente realmente busca conseguir. Los lentes inteligentes son lo bastante visibles como para que las grandes empresas no los pasen por alto mucho tiempo; Rokid debe tomar suficiente ventaja en producto, usuarios, desarrolladores y cadena de suministro antes de que las grandes empresas entren de lleno.
\end{knowledgebox}

\begin{knowledgebox}{Términos clave: los cuatro «no»}
\noindent\begin{tabularx}{\textwidth}{p{0.16\textwidth}p{0.36\textwidth}X}
\toprule
Etapa & Situación del gigante & Ventana de la empresa emergente \\
\midrule
No lo ven & La oportunidad es demasiado marginal y la organización no la ha notado & Un equipo pequeño puede validar primero una necesidad nueva. \\
No lo entienden & Vieron el fenómeno, pero no pueden juzgar su valor & La empresa emergente demuestra su relevancia con producto y datos. \\
No les interesa & Lo entendieron, pero por ahora la escala no es suficiente & Es la ventana más probable hoy para los lentes inteligentes. \\
No llegan a tiempo & Cuando el gigante entra, el pionero ya acumuló experiencia y ecosistema & La empresa emergente debe buscar llegar a este estado. \\
\bottomrule
\end{tabularx}
\end{knowledgebox}

\subsection{Por qué las grandes empresas llegan un paso atrás}

Zhu Mingming considera que las grandes empresas suelen atreverse a ir detrás de los demás: entran con cadenas de suministro maduras, soluciones maduras y menor riesgo. Por ejemplo, los lentes sin pantalla son más fáciles de fabricar con la cadena de suministro existente; en cambio, los lentes inteligentes con pantalla, ligeros y all-in-one requieren explorar procesos de fabricación nuevos, y ahí la empresa emergente puede ponerse al frente. El problema es que ir adelante implica asumir, en nombre de toda la industria, los riesgos de los procesos de fabricación, de la capacidad de producción y de la educación del usuario.

\begin{warningbox}{Ser pionero no es solo una ventaja}
El pionero debe asumir procesos de fabricación no validados, el aumento gradual de la cadena de suministro, defectos de producto que se magnifican, la gestión de las expectativas del usuario y la presión financiera. Las grandes empresas, al entrar un paso después, pueden copiar la ruta ya madura; la empresa emergente tiene que convertir esa diferencia de tiempo en reputación, desarrolladores y experiencia de producto.
\end{warningbox}

\subsection{Competencia entre pares: no es competencia en igualdad de condiciones}

Esta subsección explica lo que Zhu Mingming llama «sentarse a la mesa de los adultos». Una empresa emergente no puede competir en igualdad de condiciones con Tencent (腾讯), Alibaba (阿里), ByteDance (字节), Xiaomi (小米), Apple (苹果) o Meta, porque el capital, el tráfico, la cadena de suministro y los canales no están en el mismo orden de magnitud. Lo que sí puede conseguir es una competencia entre pares: que, en una definición de producto concreta, los usuarios estén dispuestos a comparar Rokid con los gigantes dentro del mismo conjunto de opciones; que los canales estén dispuestos a venderlo en serio; que los desarrolladores estén dispuestos a adaptar sus aplicaciones, y que los socios de modelos y de contenido estén dispuestos a colaborar.

En la entrevista se mencionan varias señales para medirlo: uso diario promedio de más de dos horas, más de diez mil desarrolladores, más de tres mil desarrolladores empresariales, actividad semanal superior al setenta por ciento y una proporción elevada de recompra y recomendación. Estas cifras no deben tomarse como conclusiones auditadas externamente, sino como el tipo de evidencia que el ponente usa para mostrar que «se están formando atributos de plataforma». Lo que de verdad importa es que una empresa emergente no puede sentarse a la mesa solo con el entusiasmo de un evento de lanzamiento: tiene que demostrar tiempo de uso, recompra, desarrolladores y capacidad de entrega de la cadena de suministro.

\noindent\begin{tabularx}{\textwidth}{p{0.20\textwidth}p{0.32\textwidth}X}
\toprule
Indicador de competencia entre pares & Por qué importa & Limitaciones \\
\midrule
Tiempo de uso & Indica que el dispositivo no es una novedad pasajera & También hay que observar la estructura de tareas y la retención. \\
Recompra/recomendación & Indica que los primeros usuarios difunden el producto de boca en boca & La muestra puede estar sesgada hacia entusiastas de la tecnología; hay que ampliarla al usuario común. \\
Número de desarrolladores & Indica los primeros indicios de una oferta del ecosistema & Aún hay que validar la actividad y los ingresos de los desarrolladores. \\
Reabastecimiento de los canales & Indica que el producto se sostiene en ventas reales & La primera distribución no equivale a canales sólidos. \\
Capacidad de producción & Indica que el entusiasmo puede convertirse en entregas & Una capacidad alta también magnifica los problemas de calidad. \\
\bottomrule
\end{tabularx}

\subsection{Ecosistema y aliados}

Rokid no cree que pueda construir sola un ecosistema completo. Zhu Mingming menciona que Tencent tiene contenido, redes sociales y escenarios de uso, y que también pueden integrarse las capacidades de modelos como Doubao (豆包), Tongyi Qianwen (通义千问) y DeepSeek. Una empresa emergente de hardware debe consolidar la plataforma, la tecnología y la experiencia de producto, pero la construcción del ecosistema requiere más aliados. El juicio realista aquí es que, si los lentes inteligentes han de convertirse en un punto de acceso, no pueden ganar solo con especificaciones de hardware: también necesitan la participación conjunta de contenido, aplicaciones, desarrolladores, modelos y canales.

\begin{importantbox}{El ecosistema no lo construye una sola empresa}
Rokid puede asumir primero la exploración y el liderazgo, y consolidar la tecnología, la plataforma y la experiencia de producto; pero los socios de contenido, redes sociales, modelos, desarrolladores y escenarios de uso tienen que entrar en conjunto. Para un punto de acceso de IA portátil, si una empresa de hardware quiere hacer a la vez el hardware, el sistema operativo, los modelos, el contenido, las aplicaciones y los canales, el consumo de recursos supera rápidamente lo que una empresa emergente puede soportar.
\end{importantbox}

\subsection{Resumen de la sección}

Los lentes inteligentes son un campo de batalla al que los gigantes sin duda entrarán. La ventana de la empresa emergente está en asumir antes los riesgos, formar más rápido la experiencia de producto y su posición en el ecosistema, y ganarse el derecho a competir entre pares antes de que los gigantes entren de lleno.
\section{El bosque oscuro del emprendimiento de hardware: cadena de suministro, capacidad de producción y organización}

Esta sección vuelve al «bosque oscuro del hardware». Lo difícil de emprender en hardware no se limita a definir el producto: también incluye la cadena de suministro, el inventario, la capacidad de producción, el rendimiento de fabricación, la entrega, los canales, el servicio posventa y el flujo de efectivo. Un producto de software puede iterarse con rapidez; en hardware, cada iteración tiene costos de materiales, procesos, moldes e inventario. Zhu Mingming (祝铭明) subrayó varias veces que, en sus primeros años, Rokid pagó el precio de aprender la cadena de suministro con hardware de IA y bocinas inteligentes, y que gracias a eso, cuando pasó a la AR, ya contaba con un equipo de hardware y una base de producción.

\begin{figure}[H]
\centering
\includegraphics[width=0.92\textwidth]{hardware-black-forest.png}
\caption{El bosque oscuro del emprendimiento de hardware: la cadena de suministro, el flujo de efectivo, el inventario, los canales y la competencia de los gigantes coexisten. Diagrama conceptual propio, elaborado a partir de la conversación de 00:59:17--02:08:56.}
\end{figure}

\begin{knowledgebox}{Lectura de la figura: el bosque oscuro es un campo de riesgos con muchas variables}
La cadena de suministro determina si se puede entregar; el flujo de efectivo, si se puede resistir hasta la siguiente generación de producto; el inventario, la exposición al riesgo; los canales, si se puede llegar a los usuarios; los gigantes, la intensidad de la competencia; y el PMF, si todo esto vale la pena. Quien emprende en hardware tiene que gestionar todas estas variables al mismo tiempo.
\end{knowledgebox}

\subsection{El precio de aprender hardware: del producto para entusiastas al producto de consumo}

Esta subsección vuelve a lo que le costaron a Rokid, en sus inicios, Alien y las bocinas. Zhu Mingming afirma que un buen producto a ojos de los entusiastas de la tecnología no es necesariamente uno que los consumidores quieran comprar. Un producto puede ser muy atractivo, muy completo y de gran complejidad técnica, y aun así no llegar a ser un producto de consumo por su precio, su peso, su escenario de uso, sus canales o la complejidad de su servicio. Lo implacable del emprendimiento de hardware es que los usuarios no pagan por la «dificultad técnica», y la cadena de suministro tampoco reduce sus costos porque la visión sea ambiciosa.

\begin{warningbox}{Iterar rápido en hardware es desastroso}
El software puede publicar versiones de forma continua; en hardware, cada cambio puede afectar moldes, materiales, líneas de producción, inventario y servicio posventa. Una empresa de hardware en etapa temprana tiene que controlar el ritmo de iteración: si es demasiado lento, el mercado la deja atrás; si es demasiado rápido, desestabiliza la cadena de suministro, el flujo de efectivo y el sistema de calidad.
\end{warningbox}

\noindent\begin{tabularx}{\textwidth}{p{0.20\textwidth}p{0.32\textwidth}X}
\toprule
Riesgo & Cómo se manifiesta en un producto de software & Efecto amplificado en un producto de hardware \\
\midrule
Iteración de versiones & Publicar parches o actualizaciones graduales & Involucra materiales, moldes, certificaciones e inventario. \\
Retroalimentación de usuarios & Pruebas A/B rápidas & Ciclo de retroalimentación largo; devoluciones, cambios y posventa costosos. \\
Defectos de calidad & Revertir o corregir en producción & Los equipos ya entregados generan indemnizaciones reales y daño a la reputación. \\
Aumento repentino de la demanda & Agregar servidores o limitar el tráfico & Exige capacidad de producción, materiales, rendimiento de fabricación y adelantar efectivo. \\
Expansión de canales & Distribución en línea de costo relativamente bajo & Exige surtir puntos de venta, capacitar, exhibir y verificar el reabastecimiento. \\
\bottomrule
\end{tabularx}

\subsection{Capacidad de producción y el efecto adverso de volverse viral}

Después de que Rokid recibiera más atención en fechas recientes, Zhu Mingming insistió en que la presión aumentó en lugar de disminuir. La popularidad trae atención para la marca, pedidos y talento, pero también le quita al producto el margen para pulirse en un círculo reducido. Antes se podía entregar un producto de 70 puntos a un grupo pequeño de entusiastas y mejorarlo con ellos hasta los 90; ahora, en cuanto se expone, cientos de millones de personas lo examinan con lupa, y unas pocas voces negativas también se amplifican. Por eso hay que reajustar el ritmo de lanzamientos, la preparación de la capacidad de producción y la madurez del producto.

\begin{importantbox}{Un producto de hardware no puede perseguir solo la atención}
La atención eleva la demanda, pero también eleva la exigencia y reduce la tolerancia a errores. Cuando un producto de hardware se entrega a gran escala, la calidad, la capacidad de producción, el servicio posventa y la reputación se ponen a prueba al mismo tiempo. Volverse viral no es el final, sino el comienzo de la prueba de estrés de la cadena de suministro.
\end{importantbox}

\subsection{El know-how del OS: las diferencias de hardware terminan reflejándose en la experiencia}

Esta subsección completa lo que Rokid considera su capacidad central: el OS. Zhu Mingming sostiene que la diferencia entre unos lentes y otros no está solo en las especificaciones del hardware, sino en el refinamiento continuo que hace el OS del consumo de energía, la conectividad, la respuesta de baja latencia y la experiencia de IA. Por ejemplo, la traducción tiene que responder en muy poco tiempo, la autonomía de la batería tiene que superar de forma perceptible a la de productos similares, y la conexión Bluetooth y entre varios dispositivos tiene que ser estable. Parecen detalles, pero en un dispositivo always-on los detalles son la condición para convertirse en la puerta de entrada.

El hardware atraviesa mesetas industriales: en ciertas etapas, la óptica, los chips, las baterías y el peso difícilmente dan saltos continuos, y otros fabricantes pueden alcanzar al líder poco a poco. En cambio, el software, la experiencia y el ecosistema pueden acumularse de manera sostenida. Este juicio explica por qué Rokid insiste siempre en el OS y en los desarrolladores, en lugar de destacar solo un lente, una cámara o el nombre de un modelo.

\begin{knowledgebox}{Términos clave: qué significa el OS en unos lentes inteligentes}
\noindent\begin{tabularx}{\textwidth}{p{0.18\textwidth}p{0.35\textwidth}X}
\toprule
Capacidad & Significado concreto & Por qué influye en la puerta de entrada \\
\midrule
Control del consumo de energía & Planificación del sistema, activación de sensores, gestión de la pantalla y del audio & Un dispositivo always-on tiene que lograr que el usuario quiera llevarlo puesto más de medio día. \\
Respuesta de baja latencia & Voz, traducción, avisos y resultados visuales que regresan con rapidez & Una latencia alta hace que el usuario vuelva al teléfono. \\
Conexión estable & Coordinación estable entre Bluetooth, teléfono, modelos en la nube y componentes locales & Un dispositivo portátil no puede desconectarse ni reconectarse con frecuencia. \\
Orquestación de la experiencia & Organizar cámara, pantalla, voz, control táctil y modelos en rutas cortas & En el fondo, competir por la puerta de entrada es competir por el costo de la ruta. \\
\bottomrule
\end{tabularx}
\end{knowledgebox}

\subsection{Espíritu lúdico y trouble maker}

El primer valor de Rokid es el «espíritu lúdico». Aquí, el espíritu lúdico no significa relajamiento, sino buscarle problemas a los detalles del producto de manera constante. Zhu Mingming se describe como un trouble maker: pregunta sin cesar si esta experiencia todavía puede mejorar o si este detalle está mal. Emprender en hardware requiere esta actitud de cuestionamiento, porque muchas diferencias no están en las diapositivas, sino en los detalles del uso, el peso, la pantalla, el sonido, la línea de producción, el inventario y la retroalimentación de los usuarios.

\begin{figure}[H]
\centering
\includegraphics[width=0.88\textwidth]{playfulness-culture.png}
\caption{La cultura del espíritu lúdico: emprender en hardware requiere una actitud de cuestionamiento al estilo trouble maker. Diagrama conceptual propio, elaborado a partir de la conversación de 01:41:35--01:48:32.}
\end{figure}

\begin{knowledgebox}{Lectura de la figura: el espíritu lúdico no es frivolidad, es no conformarse nunca con el producto}
En la figura, la curiosidad, la práctica, la disposición a arriesgar, la estética y el equipo conforman una cultura de producto. Si una empresa de hardware depende solo de procesos, tiende a hacer productos aceptables pero sin impacto; si solo tiene creatividad, tiende a no poder entregar. El espíritu lúdico tiene que ir unido a la disciplina de ingeniería.
\end{knowledgebox}

\begin{knowledgebox}{Nota de clase: el espíritu lúdico debe permitir que el equipo contradiga al fundador}
Un detalle clave de la entrevista es que Zhu Mingming llegó a rechazar los dos productos más rentables de Rokid; al final, el equipo los desarrolló y cambió su juicio con la experiencia y los datos. Esto muestra que el «espíritu lúdico» no es la creatividad de un solo jefe, sino una organización que permite a los equipos de producto, negocio e ingeniería cuestionar las decisiones con prototipos físicos.
\end{knowledgebox}

\subsection{La definición interna del éxito y del fracaso}

Esta subsección aborda los valores del cierre de la entrevista. Para Zhu Mingming, el éxito de su segundo emprendimiento no es simplemente conseguir financiamiento, salir a bolsa o volverse famoso, sino que la empresa pueda seguir el camino en el que cree y que los usuarios aprecien el producto y lo usen durante mucho tiempo. Si se aparta de lo que se propuso hacer, aunque gane dinero o salga a bolsa, no lo considera un éxito. El fracaso, en cambio, es verse obligado a abandonar la dirección en la que cree o hacer lo que no quiere hacer por objetivos de corto plazo.

Esto es especialmente importante para el emprendimiento de hardware. El camino del hardware es demasiado largo y las tentaciones de corto plazo son muchas: vender la empresa, pasar a productos más fáciles, perseguir modas, sacrificar la experiencia a cambio de escala o sustituir la reputación del producto con la IP del fundador. Zhu Mingming, por el contrario, prefiere que la popularidad del producto y de la empresa supere la suya personal, porque la IP del fundador trae una carga social y además desvía la atención de la empresa del producto hacia la persona.

\begin{importantbox}{La reputación del producto va antes que la IP del fundador}
El final de esta entrevista no es un discurso motivacional, sino un problema de gobierno corporativo: cuando una empresa obtiene atención porque su fundador se vuelve conocido fuera de su nicho, ¿cómo asegurar que esa atención termine al servicio del producto, en lugar de que la organización gire en torno a la actuación de una persona? Para una empresa de hardware, la reputación de largo plazo sigue viniendo del uso cotidiano, la autonomía, la estabilidad, el servicio posventa y el uso real, no de la visibilidad del fundador.
\end{importantbox}

\subsection{Resumen de la sección}

El núcleo del bosque oscuro del emprendimiento de hardware es que cada decisión se convierte en un costo real. Emprender en lentes inteligentes exige avanzar al mismo tiempo en la definición del producto, la cadena de suministro, la capacidad de producción, el ecosistema, el financiamiento y la cultura organizacional, y cualquiera de estos frentes puede frenar el proceso de convertirse en plataforma.
\section{Síntesis y ampliación}

Esta sección condensa todo el episodio en cinco juicios. Primero, el eje de Rokid no es un producto aislado, sino la búsqueda continua de la siguiente puerta de entrada a la computación personal. Segundo, el altavoz con IA no llegó a ser una plataforma porque su entrada, su salida y sus escenarios de uso eran demasiado limitados; la oportunidad de las gafas inteligentes como puerta de entrada proviene de su funcionamiento always-on y de su capacidad de mostrar información. Tercero, el paso de la IA a la AR en 2019 fue una continuidad estratégica, no un cambio repentino de rumbo. Cuarto, China y Estados Unidos definen de manera distinta las gafas inteligentes: en China, la población con miopía y el uso durante todo el día favorecen la lógica de IA+AR, mientras que en Estados Unidos pesa más la lógica de las sunglasses y del uso social. Quinto, la verdadera dificultad del emprendimiento en hardware es que el producto, la cadena de suministro, la capacidad de producción, el ecosistema y la competencia con los gigantes deben superar el umbral al mismo tiempo.

\begin{importantbox}{Ubicar EP104 en la serie de IA/internet de Zhang Xiaojun (张小珺)}
EP104 sigue una línea distinta de la de IA en el mundo físico de EP106/EP109: trata de la puerta de entrada a la «inteligencia portátil». Si la inteligencia corporizada es la IA que entra en el mundo físico exterior, las gafas inteligentes son la puerta de entrada por la que la IA se acerca a la percepción y al flujo de información de cada persona.
\end{importantbox}

\subsection{Síntesis conceptual: los tres umbrales de la inteligencia portátil}

Esta sección condensa de nuevo todo el texto en un modelo de juicio. Para que las gafas inteligentes se conviertan en la puerta de entrada a la inteligencia portátil, deben superar al menos tres umbrales. El primero es el umbral de uso: el peso, las lentes, la graduación, la apariencia, la autonomía de la batería y la sensación de mareo deben ser aptos para el uso diario. El segundo es el umbral de interacción: en tareas breves, el costo de la ruta debe ser claramente menor que con el teléfono, y la combinación de pantalla, voz y control táctil debe resultar lo bastante natural. El tercero es el umbral de ecosistema: el tiempo de uso, los desarrolladores, los socios de modelos, los escenarios de contenido y la recompra de los canales deben poder alimentarse entre sí y crecer.

\noindent\begin{tabularx}{\textwidth}{p{0.20\textwidth}p{0.34\textwidth}X}
\toprule
Umbral & Pregunta central & Respuesta que da EP104 \\
\midrule
Umbral de uso & ¿El usuario está dispuesto a llevarlas puestas todo el tiempo? & Empezar por quienes ya usan lentes con armazón, sin intentar convencer a fuerza a quienes se resisten. \\
Umbral de interacción & ¿Llega a lo que se busca más directamente que el teléfono? & El funcionamiento always-on, la pantalla, la voz y la dirección de la mirada reducen el costo de la ruta. \\
Umbral de ecosistema & ¿Puede pasar de producto a plataforma? & Lo validan en conjunto las dos horas de uso, los desarrolladores, las alianzas con modelos y el reabastecimiento de los canales. \\
\bottomrule
\end{tabularx}

\subsection{Conexión con las entrevistas anteriores y posteriores}

Esta sección vuelve a situar EP104 en esta serie de entrevistas de Zhang Xiaojun sobre IA e internet. EP106 y EP109 se centran en la inteligencia corporizada y discuten cómo los robots entran en el mundo físico mediante datos, simulación y modelos; EP104, en cambio, se centra en la inteligencia portátil y discute cómo la IA se acerca a la percepción y al flujo de información de cada persona. En ninguno de los dos casos se trata solo de tener «modelos más potentes»: la capacidad de los modelos tiene que encontrar el hardware, la forma de interacción y el ciclo cerrado de datos adecuados.

Si se ve la llegada de la IA al hardware como un mapa, el robot es el «extremo de acción»: se mueve, sujeta, manipula y presta servicios en lugar de la persona; las gafas inteligentes son el «extremo de percepción»: ven, escuchan, consultan, traducen, recuerdan y administran la atención en lugar de la persona. Ambos requieren una cadena de suministro de largo plazo y validación en escenarios reales, pero su estructura de riesgo es distinta. La dificultad del robot está en la fiabilidad de sus acciones y en el costo del mundo físico; la de las gafas, en la comodidad de uso, la eficiencia como puerta de entrada y la densidad del ecosistema.

\subsection{Takeaways clave}

\begin{enumerate}
\item Lo esencial de las gafas inteligentes no es tomar fotos ni lucir tecnología, sino la eficiencia como puerta de entrada portátil a la IA.
\item La capacidad de mostrar información cambiará la forma de interacción con la IA: la traducción, los avisos y el consumo de información pasarán de la voz secuencial a flujos visuales.
\item La ventana de oportunidad de una empresa emergente frente a los gigantes es limitada; debe construir barreras de experiencia y de ecosistema antes de que las grandes empresas se fijen en el mercado.
\item El éxito repentino de un producto de hardware eleva, a su vez, la presión sobre la calidad y la capacidad de producción; el tráfico no debe confundirse con el PMF en sí.
\item El caso de Rokid muestra que emprender en hardware requiere una ruta de largo plazo, capacidad en la cadena de suministro, paciencia del capital y una cultura de producto que busque problemas de forma continua.
\end{enumerate}

\subsection{Lecturas adicionales}

\begin{itemize}
\item Si interesa cómo entra la IA en el mundo físico, pueden compararse las entrevistas sobre robótica de EP106 con Wang He (王鹤), EP109 con Xie Chen (谢晨) y EP121 con Tan Jie (谭捷).
\item Si interesan el emprendimiento en tecnología dura y la cadena de suministro, puede compararse con EP111, la historia del emprendimiento en LiDAR de Li Yifan (李一帆).
\item Si interesan las aplicaciones de IA y las puertas de entrada de producto, pueden compararse las entrevistas de EP123 con ONE2X y de EP130 sobre metodología de productos de IA.
\end{itemize}

\end{document}
